% Figura V7: poda (ramos válidos, inviáveis e não explorados + tesoura)
\begin{tikzpicture}[
    nd/.style={circle, draw, minimum size=8mm, thick, font=\small, inner sep=0pt},
    valido/.style={nd, draw=btValido, fill=btValido!18},
    podado/.style={nd, draw=btPodado, fill=btPodado!18},
    neutro/.style={nd, draw=btNeutro, fill=btFundo},
    ed/.style={draw=btNeutro, thick},
  ]
  % nível 0
  \node[valido] (r) at (0,0) {};
  % nível 1
  \node[valido] (a) at (-2.2,-1.7) {};
  \node[podado] (b) at ( 2.2,-1.7) {};
  % nível 2
  \node[neutro] (a1) at (-3.2,-3.4) {};
  \node[valido] (a2) at (-1.2,-3.4) {};
  \node[neutro] (b1) at ( 1.2,-3.4) {};
  \node[neutro] (b2) at ( 3.2,-3.4) {};

  \draw[ed] (r) -- (a);
  \draw[draw=btPodado, very thick, dashed] (r) -- (b);
  \draw[ed] (a) -- (a1);
  \draw[ed] (a) -- (a2);
  \draw[ed, draw=btNeutro!55] (b) -- (b1);
  \draw[ed, draw=btNeutro!55] (b) -- (b2);

  % tesoura sobre a aresta podada
  \node[text=btPodado, font=\large, fill=white, inner sep=1pt] at ($(r)!0.5!(b)$) {\ding{34}};

  % legenda compacta
  \begin{scope}[shift={(-3.6,-4.6)}]
    \node[valido, minimum size=4mm] (lv) at (0,0) {};
    \node[anchor=west, font=\scriptsize, text=btTexto] at (0.3,0) {válido};
    \node[podado, minimum size=4mm] (lp) at (2.4,0) {};
    \node[anchor=west, font=\scriptsize, text=btTexto] at (2.7,0) {podado};
    \node[neutro, minimum size=4mm] (ln) at (4.9,0) {};
    \node[anchor=west, font=\scriptsize, text=btTexto] at (5.2,0) {não explorado};
  \end{scope}
\end{tikzpicture}
